\documentclass{article}
\usepackage{fuzz}
\begin{document}

\section{Basic Types}

\begin{zed}
  [NAME]
\end{zed}

\section{State}

A \emph{state schema} declares the variables a system holds
and the \emph{invariants} that must always be true.

\begin{schema}{Counter}
  count : \nat
\where
  count \leq 100
\end{schema}

The declaration \texttt{count : \nat} says the system holds
a natural number.  The predicate \texttt{count \leq 100} is
an invariant: no reachable state may violate it.

\begin{schema}{AddressBook}
  known : \power NAME \\
  address : NAME \pfun NAME
\where
  \dom address = known
\end{schema}

This schema shows a finite set \texttt{known} and a partial
function \texttt{address} whose domain must equal \texttt{known}.

\end{document}
